\einheitenkopf{Einheit 3 von 5 · Punkte und Werte}
\zweigzeile{Hier lernst du, Funktionswerte zu berechnen, Punkte zu prüfen und Nullstellen und Achsenschnittpunkte zu bestimmen · neu in diesem Jahr · P10 oft · baut auf: lineare Funktion (Einheit 2), lineare Gleichungen lösen, negative Zahlen}
\erg{20}{a) $5 \cdot 3 + 2$ \quad b) $2 \cdot 4 - 7$ \quad c) $3 \cdot (-4) - 2$ \quad d) $-2 \cdot (-3) + 5$ \quad e) $-(-6) + 8$}
\erg{21}{a) $f(4) = 11$ \quad b) $f(2) = 9$ \quad c) $f(5) = 19$ \quad d) $f(6) = 13$ \quad e) $f(-4) = -7$ \quad f) $f(2{,}5) = 9$}
\erg{22}{a) $2x + 4 = 10$, $x = 3$ \quad b) $3x - 5 = 7$, $x = 4$ \quad c) $5x + 3 = 28$, $x = 5$ \quad d) $4x - 2 = 22$, $x = 6$ \quad e) $-2x + 6 = 14$, $x = -4$ \quad f) $4x + 7 = 1$, $x = -1{,}5$}
\erg{23}{a) $x = 4$ \quad b) $x = 3$ \quad c) $x = 2$ \quad d) $x = 7$}
\erg{24}{e) $x = -3$ \quad f) $x = 3$ \quad g) $x = -1{,}5$ \quad h) $x$-Achse: $(3{,}5|0)$, $y$-Achse: $(0|7)$}
\erg{25}{a) ja ($2 + 4 = 6$) \quad b) nein ($2 \cdot 3 - 1 = 5 \neq 6$) \quad c) nein ($4 \cdot 1 - 2 = 2 \neq 3$) \quad d) ja ($-5 + 3 = -2$) \quad e) ja ($3 \cdot (-2) + 2 = -4$) \quad f) nein ($-2 \cdot (-3) + 1 = 7 \neq -5$) \quad g) $\frac{3}{4} \cdot (-8) - 2 = -6 - 2 = -8$, wahr: $V$ liegt auf dem Graphen}
\erg{26}{a) $-5;\ -4;\ -3;\ -2;\ -1$ \quad b) $7;\ 5;\ 3;\ 1;\ -1$ \quad c) $-3;\ -2;\ -1;\ 0;\ 1$ \quad d) $f(x) = 3x - 4$}
\erg{27}{a) $x = 2$ \quad b) $(0|{-4})$ \quad c) $x = 5$ \quad d) $(0|5)$ \quad e) $(3|2)$}
\erg{28}{a) Leon hat $x = 0$ eingesetzt und den Schnittpunkt mit der $y$-Achse berechnet. Für die Nullstelle setzt man $f(x) = 0$: $0 = 2x - 6$, $x = 3$, Nullstelle bei $(3|0)$. \quad b) $0 = 4x - 10$, $x = 2{,}5$; $y$-Achse: $(0|{-10})$}
\erg{29}{a) wahr \quad b) falsch (Nullstelle 2) \quad c) wahr ($-3 \cdot (-1) + 6 = 9$) \quad d) falsch (bei $(0|6)$)}
\erg{30}{a) Der Graph ist die waagerechte Gerade $y = 5$; sie trifft die $x$-Achse nie, $g(x) = 0$ hat keine Lösung. \quad b) $0 = m \cdot x + n$ hat für $m \neq 0$ genau die Lösung $x = -n : m$; eine nicht waagerechte Gerade schneidet die $x$-Achse genau einmal.}
\erg{31}{a) $y = -2{,}5 \cdot 12 + 90 = 60$, also 60\,\% \quad b) $-2{,}5x + 90 = 40$, $x = 20$, nach 20 km \quad c) $-2{,}5x + 90 = 0$, $x = 36$, nach 36 km}
